% Final integral-quotient auxiliary; combined geometric application remains separate.
% Supersedes only Step 3 of the retained Witt-lift draft, whose proof remains valid.
\paragraph{Auxiliary proposition for C.3bl (Proved):
the primitive relative Milnor invariant kills the entire exceptional image.}
Let $k$ be algebraically closed of characteristic $p>0$, let $s\geq1$
and $m\geq0$, and put $r=s+1$. Let $W_1,\ldots,W_s$ be smooth
cubic surfaces, $A_i=\mathcal O_{W_i}(1)=-K_{W_i}$, and let $C/k$
be a smooth projective connected curve of genus at least two. Set
\[
 \begin{gathered}
 T=\mathbb P^m\times\prod_{i=1}^sW_i,
 \qquad L_T=\mathcal O_{\mathbb P^m}(1)\boxtimes A_1\boxtimes\cdots\boxtimes A_s,
 \qquad b=c_1(L_T),\\
 V=T\times C,\qquad L=L_T\boxtimes K_C.
 \end{gathered}
 \tag{PB.1}
\]
Suppose $R$ is an excellent henselian local $k$-algebra obtained by
strict henselization from an algebraic local $k$-scheme, and
$\mathcal Z\to\operatorname{Spec}R$ is projective and regular.
Suppose that its closed fibre is the smooth Cartier divisor
$i:V\hookrightarrow\mathcal Z$, that
\[
 \mathfrak m_R\mathcal O_{\mathcal Z}=I_V,
 \qquad I_V/I_V^2=L,
 \qquad q:\mathcal Z\longrightarrow T,
 \quad q|_V=\operatorname{pr}_T.
 \tag{PB.2}
\]
Write $E=V_2$ and $X=\mathcal Z\setminus V$. We require the given
scheme and morphism to descend to the algebraic local model, as they
do in the total-line construction of C.3bj. Every finite list of
line bundles and Chow relations used below can then be considered on
one smooth finite-type model after a pointed etale base change.

Define the following direct K\"unneth summand:
\[
 \mathcal M_k=
 H^0(\mathbb P^m,\mathcal O(1))\otimes_k
 \bigotimes_{i=1}^s H^1(W_i,A_i\otimes\Omega^1_{W_i/k})
 \otimes_k H^1(C,K_C)
 \subset H^r(V,L\otimes\Omega^s_{V/k}).
 \tag{PB.3}
\]
Let $P$ be the projection onto this summand, with the indicated order
of the cubic factors. If $\mathcal A$ is a line bundle on $\mathcal Z$
trivialized on $V$, write $a\in H^1(V,L)$ for its first relative jet.
For $D_i\in\operatorname{Pic}(W_i)$ form the actual Chow class
\[
 \beta=c_1(\mathcal A)\prod_{i=1}^s c_1(q^*D_i)
                      \in\mathrm{CH}^{r}(\mathcal Z).
 \tag{PB.4}
\]
Here $q^*D_i$ abbreviates pullback through the $i$th projection of $T$.
If
\[
 P\left(a\smile\prod_i d\log D_i\right)\ne0,
 \tag{PB.5}
\]
then $\beta|_X\ne0$ integrally. More precisely, the projected
relative Milnor invariant in (PB.5) vanishes on every exceptional
class $i_*\gamma$ whose restriction to $V$ is zero. We do not assert
that all Hodge components of every exceptional class vanish.

\begin{proof}
\textbf{1. The canonical relative class and its products.}
We use the ordinary Milnor symbol sheaves and put
\[
 \mathcal J_n=\ker(\mathcal K_n^M(E)\longrightarrow\mathcal K_n^M(V)).
\]
The separately proved relative Milnor lemma (RM.1)--(RM.7) gives
an exact sheaf sequence, $p\mathcal J_n=0$, and the canonical quotient
\[
 \tau_n:\mathcal J_n\longrightarrow L\otimes\Omega^{n-1}_{V/k},
 \qquad
 \{1+u,v_2,\ldots,v_n\}\longmapsto
 u\otimes d\log\bar v_2\wedge\cdots\wedge d\log\bar v_n.
 \tag{PB.6}
\]
This includes $p=2$. It is obtained from the rank-one conormal
kernel in the exterior powers of $\Omega^1_{E/k}|_V$, and uses no
Milnor--Quillen comparison on $E$.

The separately proved Rost-localization induction (K2.4)--(K2.6)
gives divisibility of $H^{n-1}(V,\mathcal K_n^M)$ for every $n$.
Its exact primary inputs are
\citet[Section 5, Corollary 6.5, Proposition 8.6, and Section 12,
pp.~356--360, 370, 382--383]{cyRost1996} and the integral Gersten
theorem of \citet[Theorem 7.1, author pp.~21--22]{cyKerz2009Milnor}.
Consequently
\[
 H^n(V,\mathcal J_n)\lhook\joinrel\longrightarrow
 H^n(E,\mathcal K_n^M)
 \tag{PB.7}
\]
is injective, with image the kernel of reduction to
$H^n(V,\mathcal K_n^M)$. Indeed the preceding connecting map has
divisible source and $p$-killed target. Thus every absolute class
in this kernel has a unique relative lift.

Multiplication of symbols gives
$\mathcal J_q\otimes\mathcal K_{n-q}^M(E)\to\mathcal J_n$.
After applying (PB.6), this product becomes the product of
$\tau_q$ with the reduced $d\log$ class of the second factor.
This identity is checked on the displayed symbols, in their fixed
order, and then passes to sheaf cohomology. In particular a product
of one nil line bundle with $n-1$ actual line bundles has relative
image exactly its first jet times their reduced divisor classes.
Distinct ordered cubic factors contribute with coefficient one;
there is no factorial to invert.

\textbf{2. Removing the free part of the exceptional kernel.}
Write $B_j=\mathrm{CH}^j(T)$. The integral projective-bundle and
point-blow-up formulas make each $B_j$ finite free and give
\[
 \mathrm{CH}^{s}(V)=B_s\oplus
                 \bigl(B_{s-1}\otimes_{\mathbb Z}\operatorname{Pic}(C)\bigr).
 \tag{PB.8}
\]
The ring is generated by divisors, and these isomorphisms retain
all torsion of $\operatorname{Pic}(C)$; they follow by applying the
formulas to each cubic blow-up after product with the other factors.
The exact inputs are
\citet[Theorems 2.4 and 2.8(3)]{cyJiang2023},
\citet[Lemma 42.36.2, Tag 02TX]{cyStacksProject}, and the
all-characteristic cubic blow-up description
\citep[Chapter 4, Proposition 2.4, p.~178]{D_HuybrechtsCubics}.

We recall the complete elementary injectivity argument, so that no
integral hard Lefschetz assertion is hidden here. On a cubic, over
$\mathbb Q$, multiplication by $A$ has the chain
$1\mapsto A\mapsto3z\mapsto0$ and kills $A^\perp$ in degree one.
Give this chain squared norms $1,2,4$ and define the adjoint lowering
operator by $F(A)=2$, $F(3z)=2A$. It is zero on $1$ and on $A^\perp$,
which is given any positive inner product. On $\mathbb P^m$ take
$F(h^j)=j(m-j+1)h^{j-1}$ with the corresponding positive norms.
Their tensor product has multiplication $E_b$ by $b$, adjoint $F$,
and
\[
 [E_b,F]|_{B_j\otimes\mathbb R}
       =(2j-2s-m)\operatorname{id}.
\]
For a vector of weight $w$ we have
$\|E_bv\|^2=\|Fv\|^2-w\|v\|^2$.
Thus $E_b$ is injective in negative weight, and $E_b^2$ is injective
in weight at most $-2$. At weight $-2$, if $E_b^2v=0$, put
$u=E_bv$; the weight-zero identity gives $Fu=0$, and then
$\|u\|^2=\langle v,Fu\rangle=0$. In degree $s-1$ the weight is
$-m-2$. Hence both $b$ and $b^2$ are injective on $B_{s-1}$,
including $m=0$, since this integral group is free.

Choose $o\in C(k)$, and write $c_1(K_C)=e[o]+\kappa_0$, where
$e=2g(C)-2\ne0$ and $\kappa_0\in\operatorname{Pic}^0(C)(k)$.
For $\gamma=\alpha+v[o]+\eta$ according to (PB.8), the equation
$c_1(L)\gamma=0$ gives
\[
 b\alpha=0,\qquad e\alpha+bv=0,
 \qquad
 \eta\in B_{s-1}\otimes\operatorname{Pic}^0(C)(k).
\]
It follows that $b^2v=0$, hence $v=0$ and then $\alpha=0$.
The remaining equation is exactly $b\eta=0$.
Choose integral Smith bases for the injective lattice map
$b:B_{s-1}\to B_s$, writing
\[
 bv_j=d_jw_j,\quad d_j>0,
 \qquad
 \gamma=\sum_jv_jQ_j,\quad
 Q_j\in\operatorname{Pic}^0(C)(k),\quad d_jQ_j=0.
 \tag{PB.9}
\]
This is an exact description of the whole kernel. Nonunit $d_j$
are permitted and are not replaced by units modulo $p$.

\textbf{3. A free integral primitive quotient kills the saturated image.}
For $z\in B_s$ define
\[
 \Phi_k(z):H^0(T,L_T)\longrightarrow
 H^0(\mathbb P^m,\mathcal O(1))\otimes
            \bigotimes_i H^1(W_i,A_i\otimes\Omega^1_{W_i/k})
 \tag{PB.10}
\]
by multiplying a section by the divisor-cycle $d\log$ class of $z$
and projecting to the indicated component. Products of divisor
classes realize this map; rationally equivalent divisor polynomials
give the same class by the integral product compatibility
\citep[Section 3, diagram (22), pp.~557--558]{cyBraunling2013}.
In the multigrading of the integral ring of $T$, only the part
\[
       (\text{degree on }\mathbb P^m;\text{cubic degrees})
                         =(0;1,\ldots,1)
 \tag{PB.11}
\]
can contribute. On $\delta_1\boxtimes\cdots\boxtimes\delta_s$
in that part the map is the tensor product of
\[
 \mu_{\delta_i}:H^0(W_i,A_i)\longrightarrow
           H^1(W_i,A_i\otimes\Omega^1_{W_i/k}),
 \qquad u\longmapsto u\,d\log\delta_i,
\]
and the identity on the projective-space factor.

Each $A_i$ is primitive in the free Picard lattice: its intersection
number with a line is one. Thus
\[
 \Lambda_i=\operatorname{Pic}(W_i)/\mathbb ZA_i,
 \qquad \Lambda=\bigotimes_{i=1}^s\Lambda_i
\]
are free abelian groups. Let $\Psi:B_s\to\Lambda$ be the integral
projection to the multidegree in (PB.11), followed by the quotient
in each cubic Picard factor. For $v\in B_{s-1}$, the contribution
of $hv$ to that multidegree is zero. Every other term of $bv$ that
reaches it inserts an $A_i$ into a missing cubic factor. Therefore
\[
                         \Psi(bv)=0.
\]

We verify that $\Phi_k$ factors through $\Psi$ in every characteristic.
The restricted Euler sequence and cubic hypersurface cohomology give
\[
 \begin{gathered}
 0\longrightarrow\Omega^1_{\mathbb P^3/k}(1)|_{W_i}
       \longrightarrow\mathcal O_{W_i}^{\oplus4}
       \longrightarrow A_i\longrightarrow0,\\
 H^1(W_i,\Omega^1_{\mathbb P^3/k}(1)|_{W_i})=0.
 \end{gathered}
 \tag{PB.12}
\]
Indeed the global middle-to-right map is the isomorphism from the
four coordinates to $H^0(W_i,A_i)$, and
$H^1(W_i,\mathcal O_{W_i})=0$. These cohomology statements follow
from the cubic hypersurface sequence and
\citet[Lemmas 30.8.1--30.8.2, Tag 01XS]{frPrimitiveProjectiveCohomology}.
The class $d\log A_i$ comes from the ambient cotangent sheaf, so
every section times that class factors through the zero group in
(PB.12). Hence $\mu_{A_i}=0$. The maps $\mu_{\delta_i}$ are additive
in the integral Picard class and therefore factor through $\Lambda_i$.
Their tensor product proves the desired factorization of $\Phi_k$.

Apply $\Psi$ to the Smith relations (PB.9). Since $\Lambda$ is free,
\[
 d_j\Psi(w_j)=\Psi(bv_j)=0
 \quad\Longrightarrow\quad\Psi(w_j)=0
 \quad\Longrightarrow\quad\Phi_k(w_j)=0.
 \tag{PB.13}
\]
The first implication is in a free integral lattice before passing
to characteristic $p$. It remains valid when $p\mid d_j$ and proves
that $\Phi_k$ kills the entire saturation of $bB_{s-1}$.
This argument needs neither a division by a Smith invariant factor
nor a cohomological hard Lefschetz theorem.

\textbf{4. Actual lifts of curve line bundles and a divisible-group vanishing.}
The product cohomology calculations give
\[
 \begin{split}
 H^1(V,L)&=H^0(T,L_T)\otimes H^1(C,K_C),\\
 H^2(V,L^n)&=0\quad(n\ge1),\qquad
 H^1(V,L^n)=0\quad(n\ge2).
 \end{split}
 \tag{PB.14}
\]
Every line bundle pulled from $C$ consequently lifts through all
neighbourhoods of $V$ by the square-zero unit sequence
\citep[Lemma 37.4.1, Tag 0C6R]{frC3azStacksPicThickening}. The ideals of those neighbourhoods agree with
the powers of the base maximal ideal by (PB.2). Formal existence
over $\widehat R$ algebraizes the compatible invertible system
\citep[Lemma 30.24.3, Tag 0885]{frC3azStacksExistence}.
It is invertible everywhere because it is invertible on the closed
fibre and its noninvertible locus has proper image avoiding the
closed point. Popescu's theorem for the regular completion map of
the excellent ring, finite-presentation descent of a line bundle
and its inverse, and smooth henselian lifting preserving its first
jet then produce an actual line bundle on $\mathcal Z$.
These are exactly
\citet[Lemma 15.51.14, Tag 0AH2]{frC3azStacksCompletion},
\citet[Theorem 16.12.1, Tag 07GC]{frC3azStacksPopescu},
\citet[Lemma 32.10.3(2), Tag 0B8W]{frC3azStacksLimitBundles}, and
\citet[Lemma 15.13.3, Tag 0H74]{frC3azStacksHenselSmooth}.
Only finitely many individual line bundles are needed here; no
unproved simultaneous algebraization of an entire Picard functor
is used.

There is a canonical actual extension of $K_C$ for this purpose:
\[
 \widetilde K=\mathcal O_{\mathcal Z}(-V)\otimes q^*L_T^{-1},
 \qquad \widetilde K|_V=\operatorname{pr}_C^*K_C.
 \tag{PB.15}
\]
In particular we do not assume a projection $\mathcal Z\to C$.
For any $Q\in\operatorname{Pic}^0(C)(k)$ choose an actual lift
$\widetilde Q$ on $\mathcal Z$. The product
$c_1(\widetilde K|_E)c_1(\widetilde Q|_E)$ reduces to zero in
$H^2(V,\mathcal K_2^M)$, since both reduced classes are pulled from
a curve and their Chow product there is zero. By (PB.7) it has a
unique relative lift $\xi_Q\in H^2(V,\mathcal J_2)$.
Define
\[
 F(Q)=H^2(\tau_2)(\xi_Q)
                      \in H^2(V,L\otimes\Omega^1_{V/k}).
 \tag{PB.16}
\]
This is independent of the chosen lift. Indeed two choices differ
on $E$ by a line bundle trivial on $V$, whose first jet is an
element $a'\in H^1(V,L)$. The relative product formula in Step 1
changes (PB.16) by $a'\smile d\log K_C$. By (PB.14) both factors
have cohomological degree one on $C$, so the result belongs to
$H^2(C,K_C\otimes\Omega_C^1)=0$ and is zero.

We may use tensor products to lift a sum of two $Q$'s. Uniqueness
in (PB.7) therefore proves that $F$ is an additive homomorphism.
Its source $\operatorname{Pic}^0(C)(k)$ is divisible, whereas its
target is a vector space in characteristic $p$. Hence
\[
                         F(Q)=0\quad\text{for every }Q.
 \tag{PB.17}
\]
This argument proves vanishing of the indicated relative differential
class; it does not claim that the relative Milnor class $\xi_Q$
itself is zero.

\textbf{5. The exact exceptional Chow identity and its invariant.}
Suppose, for contradiction, that $\beta|_X=0$. Chow localization
gives $\beta=i_*\gamma$. Since $\mathcal A|_V$ is trivial,
self-intersection gives $c_1(L)\gamma=0$. Apply (PB.9), choose
actual lifts $\widetilde Q_j$, and pull all $v_j,w_j$ through $q$.
The projection formula and the actual line-bundle identity (PB.15)
give the following identity in $\mathrm{CH}^r(\mathcal Z)$:
\[
 \begin{split}
 i_*\gamma
  &= -\sum_j(b+c_1(\widetilde K))v_jc_1(\widetilde Q_j)\\
  &= -\sum_j\left(
          w_jc_1(\widetilde Q_j^{\,d_j})
          +v_jc_1(\widetilde K)c_1(\widetilde Q_j)\right).
 \end{split}
 \tag{PB.18}
\]
Every class in this equation is an integral polynomial in actual
line bundles; $v_j,w_j$ are realized by divisor polynomials on the
cellular variety $T$. The equation uses neither a chosen rational
division by $d_j$ nor an equality merely in cohomology.

For the first terms, $\widetilde Q_j^{\,d_j}|_V$ is trivial.
Their relative Milnor lifts therefore have differential equal to
their first jet times the divisor-cycle class of $w_j$.
The projection $P$ of this class is zero by (PB.13) and (PB.14).
For the second terms use the relative class $\xi_{Q_j}$ from
Step 4 and multiply it by the actual class of $v_j$.
Its image under $\tau_r$ is zero by (PB.17) and the symbol-product
identity of Step 1. Each term of (PB.18) thus reduces to zero on
$V$ and has the described relative lift. Their sum is the unique
relative lift of their absolute sum by (PB.7).

For precision in passing from (PB.18) to Milnor cohomology, descend
its finitely many line bundles, cycles and rational-equivalence
witnesses to one pointed etale neighbourhood of the algebraic local
model. Shrink that base around its distinguished point so that the
proper source is regular; its nonregular locus has proper image
avoiding the point. The resulting finite-type source is smooth
over the perfect field $k$. The etale map with residue field $k$
preserves the quotient by $\mathfrak m^2$, so the actual $E$ and all
the prescribed first jets remain unchanged. The integral
Milnor--Bloch formula and its ring-product compatibility on this
smooth model, not on the nonreduced $E$, identify the Chow identity
with the corresponding cup identity
\citep[Theorems 7.1 and 7.5, author pp.~21--22]{cyKerz2009Milnor}
\citep[Section 3, diagram (22), pp.~557--558]{cyBraunling2013}.
Restrict that cup identity to $E$, and then use the uniqueness of
relative lifts in (PB.7). The projected relative differential of
the right side is zero, whereas that of $\beta$ is (PB.5).
This contradiction proves the proposition.
\end{proof}

\noindent\emph{Scope.}
The auxiliary proposition holds also for $m=0$ under its stated
geometric hypotheses. It does not, by itself, construct those
hypotheses, an ordinary algebraic family, a nonzero cup in (PB.5),
or a principal target on which the actual cycle is trivial.
Those assertions belong to a separately audited combined local
construction. The proof handles the entire exceptional image and
retains every Smith invariant factor; its claim is stronger than
checking only the numerical class or the nonsaturated $b$-image.


\noindent\emph{Proof history.}
The complete earlier draft proves (PB.13) by lifting the designated
Picard lattice to the Witt ring and using a free twisted cotangent
cohomology module. That independent proof is retained unchanged.
The present Step 3 replaces it by the free integral quotient
$\bigotimes_i(\operatorname{Pic}(W_i)/\mathbb ZA_i)$ and therefore
removes the lifting argument from this final proof. No assertion
of the earlier proof has been withdrawn.
